\begin{longtable}{@{}P{.07\linewidth}P{.47\linewidth}P{.23\linewidth}P{.15\linewidth}@{}}
\caption{Aanvullende traceability: eigen inferentiestap, alternatief en toekomstige experimenten. Alle bronlocaties staan in het derivation-register.}\\
\toprule Eis & Aanvullende ontwerpaanname & Concurrerend alternatief & Propositie / test \\ \midrule\endfirsthead
\toprule Eis & Aanvullende ontwerpaanname & Concurrerend alternatief & Propositie / test \\ \midrule\endhead
R01 & Rollen moeten voor beheerders en afnemers onderscheidbaar zijn; dat vereist geen universele resource-disjunctie. & Gedeelde resource met gedocumenteerde meervoudige typering & TP1\newline E1 \\
R02 & Een gezagsvraag kan alleen worden geadjudiceerd als definitie en context naar een controleerbare bron verwijzen. & Conventionele bronverwijzing en handmatige dossiercontrole & TP1, TP2\newline E1, E2 \\
R03 & Persistentie is een beleidseigenschap die niet door URI-syntax alleen wordt gerealiseerd. & Lokale sleutel met stabiele scope-/resolutietabel & TP1, TP4\newline E1, E4 \\
R04 & De gegevenselementmapping voegt interpretatie-informatie toe; het gebruik van een 11179-register is niet uniek noodzakelijk. & Goed gedocumenteerd conventioneel datadictionary & TP1\newline E1 \\
R05 & Decodeerbaarheid wordt vereist als contractdoel; ISO-scope geldt uitsluitend voor relevante karakteristieke masterdataberichten. & Andere syntax met gelijkwaardige dataspecificatie & TP1\newline E1 \\
R06 & Historische taken vragen twee expliciete tijdsperspectieven; OWL-Time is slechts een representatiekandidaat. & Bitemporele relationele tabellen & TP4\newline E4 \\
R07 & Afleidingsketens maken beoordeling mogelijk, maar provenance kan onjuist zijn en bewijst geen gezag. & Versiebeheerd log met dezelfde bron-/regelgegevens & TP1, TP2, TP4\newline E1, E2, E4 \\
R08 & Authenticiteit, betrouwbaarheid en sender trust zijn verschillende beslissingen; hun scheiding is eigen profielbeleid. & Brongebonden gezagsbesluit in conventionele metadata & TP2\newline E2 \\
R09 & Een score zonder populatie en methode is niet vergelijkbaar; de gekozen verplichte contextvelden zijn ontwerpkeuze. & Kwaliteitsrapport met dezelfde context buiten de graaf & TP1\newline E1 \\
R10 & Versieafhankelijkheden kunnen impactanalyse ondersteunen; meer afhankelijkheden kunnen onderhoudslast verhogen. & Zelfde versiearchief met handmatig bijgehouden releaseregister & TP4\newline E4 \\
R11 & Motiveer een mapping om onterechte equivalentie te kunnen beoordelen; labels alleen zijn geen geldige goudstandaard. & Handmatig geadjudiceerde mappingtabel & TP1, TP3\newline E1, E3 \\
R12 & Verantwoordelijke statusovergangen zijn governancekeuze; BOMOS schrijft niet deze volledige operationele lifecycle voor. & Bestaande workflow met beslislog & TP4\newline E4 \\
R13 & Productcontractvelden moeten eerst op DCAT/DPROD worden gemapt; eigen metadata zijn alleen conditionele aanvulling. & Bestaand DPROD/DCAT-profiel zonder SMDP-uitbreiding & TP1\newline E1 \\
R14 & Profielconformiteit en semantische rondreis zijn aparte toetsen; een geslaagde interface voorspelt geen begrip. & OpenAPI plus conventionele contextlinks & TP1\newline E1 \\
R15 & Eventenvelop adresseert niet alle domein-/herstelbetekenis; additionele regels hangen af van de afnemertaak. & Polling met versiearchief in plaats van events & TP4\newline E4 \\
R16 & Inferentie, graafconformiteit en werkelijkheid gebruiken verschillende toetsgronden. & SQL-/schemavalidatie plus aparte externe kwaliteitscontrole & TP1\newline E1 \\
R17 & Machinebeschikbaarheid is geen gegarandeerd juist gebruik; betere toegang kan een rivaliserende verklaring zijn. & Zelfde context in conventioneel retrievalpakket & TP1, TP2\newline E1, E2 \\
R18 & FDS-afspraken zijn contextgebonden; toelatingsdekking is geen effectmetric. & Handmatige checklist met hetzelfde toepasselijkheidsbewijs & TP1\newline E1 \\
R19 & Terugmeldingen moeten routeerbaar zijn; een eigen eventstack is niet noodzakelijk. & Bestaand terugmeldproces met versie-/besluitverwijzing & TP4\newline E4 \\
R20 & Bewaarmetadata zijn alleen relevant na expliciet bewaarbeleidsbesluit; geen generieke bewaarplicht afgeleid. & Los archiefdossier gekoppeld aan release & TP4\newline E4 \\
R21 & Juridische toepasselijkheid vereist een deskundig contextbesluit; een complianceveld volstaat niet. & Gemotiveerd juridisch dossier buiten productgraaf & TP1, TP2\newline E1, E2 \\
R22 & Ontologische verdieping is alleen kandidaat bij relevante categoriefouten; nut is kostenafhankelijk. & MIM-review met dezelfde expertise en taakset & TP3\newline E3 \\
R23 & Registratie, productiemethode, onzekerheid en gezag zijn onafhankelijke profielvelden; PROV-O schrijft deze combinatie niet voor. & Conventioneel contextrecord met dezelfde annotaties & TP2\newline E2 \\
\bottomrule\end{longtable}